\chapter*{Appendix E: Topos-Theoretic Foundations of RSVP}
\addcontentsline{toc}{chapter}{Appendix E: Topos-Theoretic Foundations of RSVP}

\section{Elementary Topoi}

\subsection{Definition}

\begin{definition}[Elementary Topos]
An elementary topos $\mathscr{E}$ is a category with:
\begin{enumerate}
\item all finite limits,
\item exponentials: for all $A,B$, an object $B^A$ with evaluation map,
\item a subobject classifier $\Omega$.
\end{enumerate}
\end{definition}

\section{Internal Logic}

\subsection{Heyting Structure}

The subobject classifier $\Omega$ carries a Heyting algebra structure.  
Thus $\mathscr{E}$ supports intuitionistic higher-order logic internally.

\begin{proposition}
Every topos is a model of intuitionistic higher-order logic.
\end{proposition}

\section{Sheaf Topoi}

\subsection{Sheaves on a Site}

Let $(\mathcal{C},J)$ be a site.

\begin{definition}[Sheaf Topos]
The category $\mathbf{Sh}(\mathcal{C},J)$ of $J$-sheaves is a topos.
\end{definition}

Objects are sheaves, morphisms are natural transformations.

\section{RSVP and Polyxan Sites}

\subsection{Semantic Site}

Define the semantic site:
\[
\mathcal{C}_{\mathrm{RSVP}}
  = \{ U\subseteq M \ \text{open} \},
\qquad
J_{\mathrm{cov}} = \text{open covers}.
\]

\subsection{Hypergraph Site}

Define the Polyxan site:
\[
\mathcal{C}_{\mathrm{Polyx}}
  = \{\text{incidence subhypergraphs of $\mathcal{G}$}\},
\qquad
J_{\mathrm{sub}} = \text{subhypergraph covers}.
\]

\section{Topos Pair}

\[
\mathscr{E}_{\mathrm{RSVP}} = \mathbf{Sh}(M),
\qquad
\mathscr{E}_{\mathrm{Polyx}} = \mathbf{Sh}(\mathrm{Inc}(\mathcal{G})).
\]

These two topoi form the semantic–algebraic foundation of the RSVP–Polyxan universe.

\section{Geometric Morphisms}

A geometric morphism
\[
f : \mathscr{E} \to \mathscr{F}
\]
is a pair of adjoint functors:
\[
f^\ast : \mathscr{F}\to \mathscr{E}, \qquad
f_\ast : \mathscr{E}\to \mathscr{F},
\]
with $f^\ast$ left exact.

\section{Synthetic Duality as a Geometric Morphism}

Let
\[
\mathbb{D} : \mathscr{E}_{\mathrm{RSVP}} \leftrightarrows
               \mathscr{E}_{\mathrm{Polyx}} : \mathbb{D}^{-1}
\]
be the duality functor of Chapter~62.

\begin{proposition}
$\mathbb{D}$ is a geometric morphism.
\end{proposition}

\begin{proof}
$\mathbb{D}^\ast$ is defined by pulling back embeddings along Polyxan incidence; it preserves finite limits.  
Its right adjoint $\mathbb{D}_\ast$ is given by taking global sections over incidence fibers.  
Thus $(\mathbb{D}^\ast,\mathbb{D}_\ast)$ is a geometric morphism.
\end{proof}

\section{Internal RSVP Fields}

Let $\underline{\mathbb{R}}$ be the constant sheaf of real numbers.

\begin{definition}[Internal Fields]
\[
\Phi \in \mathrm{Hom}_{\mathscr{E}_{\mathrm{RSVP}}}(1,\underline{\mathbb{R}}^M),
\qquad
\mathbf{v} \in \mathrm{Hom}(1,TM),
\qquad
S \in \mathrm{Hom}(1,\underline{\mathbb{R}}^M).
\]
\end{definition}

Thus RSVP fields are global elements of exponential objects.

\section{Internal Polyxan Structures}

For the incidence object $\underline{\mathrm{Inc}(\mathcal{G})}$:

\[
\iota \in \mathrm{Hom}_{\mathscr{E}_{\mathrm{Polyx}}}
 ( \underline{\mathrm{Inc}(\mathcal{G})},
   \underline{M} )
\]

is an internal embedding.

\section{Subobject Classifier and Modalities}

Let $\Omega_{\mathrm{RSVP}}$ and $\Omega_{\mathrm{Polyx}}$ be subobject classifiers in their respective topoi.

\subsection{Modal Operators}

Define modalities:
\[
\Box_{\mathrm{RSVP}} : \Omega_{\mathrm{RSVP}} \to \Omega_{\mathrm{RSVP}},
\qquad
\Box_{\mathrm{Polyx}} : \Omega_{\mathrm{Polyx}} \to \Omega_{\mathrm{Polyx}}.
\]

\begin{proposition}
Both $\Box$ operators arise as Lawvere–Tierney topologies.
\end{proposition}

\section{Lawvere–Tierney Topologies}

A Lawvere–Tierney topology on $\mathscr{E}$ is:
\[
j : \Omega \to \Omega
\]
satisfying:
\[
\top \le j(\top), \qquad
j(p\wedge q)=j(p)\wedge j(q), \qquad
j(j(p)) = j(p).
\]

\begin{proposition}
$\Box_{\mathrm{RSVP}}$ and $\Box_{\mathrm{Polyx}}$ are idempotent modalities corresponding to harmonicity and quine-stability.
\end{proposition}

\section{Sheafification via Modalities}

For any object $A$:
\[
A_j = \{ s : 1\to A \mid j(\chi_s)=\chi_s \}
\]
is the $j$-sheafification.

\begin{proposition}
Harmonic RSVP fields are the $\Box_{\mathrm{RSVP}}$-sheafification of arbitrary fields.
\end{proposition}

Similarly for Polyxan structures.

\section{Classifying Topos}

\subsection{Definition}

Let $\mathcal{T}_{\mathrm{RSVP}}$ be the theory of RSVP fields.  
Let $\mathbf{Set}[\mathcal{T}_{\mathrm{RSVP}}]$ be its classifying topos.

\begin{proposition}
$\mathscr{E}_{\mathrm{RSVP}}$ is a model of $\mathcal{T}_{\mathrm{RSVP}}$.
\end{proposition}

\section{Topos of Curvature Strata}

Let $\mathcal{S}$ be the poset of curvature strata.  
Define the site:
\[
\mathcal{C}_{\mathrm{strata}} = \mathcal{S}^{\mathrm{op}},
\qquad
J = \text{coverage by upward-closed families}.
\]

The associated topos encodes the stratified structure central to the Assembly–Curvature Correspondence.

\section{Exponentials and Embedding Flows}

Embedding flows
\[
t \mapsto \iota_t
\]
are internal morphisms:
\[
1 \to M^T,
\]
where $T$ is the time object.

\section{Universes}

\begin{definition}[Grothendieck Universe]
A Grothendieck universe $U$ is a set closed under pairing, power set, union, and Cartesian product.
\end{definition}

We fix universes to stratify RSVP–Polyxan constructions.

\section{Topos-Theoretic Interpretation of Normal Forms}

Let $N$ be the sheaf of normalised patches.

\begin{proposition}
Normalization is a reflector:
\[
N : \mathscr{E}_{\mathrm{Polyx}} \to \mathscr{E}_{\mathrm{Polyx}}
\]
left adjoint to the inclusion of normal-form sheaves.
\end{proposition}

\section{Exactness Properties}

\begin{proposition}
Both $\mathscr{E}_{\mathrm{RSVP}}$ and $\mathscr{E}_{\mathrm{Polyx}}$ are:
\begin{enumerate}
\item extensive,
\item locally cartesian closed,
\item exact,
\item infinitary coherent.
\end{enumerate}
\end{proposition}

These properties support the existence of reconciliation, duality, and patch monads.

\section{Sheaf Semantics of Reconciliation}

Reconciliation is the left Kan extension:
\[
\mathcal{R} = \mathrm{Lan}_\iota(\mathrm{id}),
\]
computed internally in the topos of Polyxan sheaves.

\section{Internal Categories}

Polyxan incidence categories and RSVP flow groupoids are internal categories in their respective topoi.

\section{Topos-Theoretic Synthetic Duality}

\begin{proposition}
Synthetic duality is an equivalence of topoi:
\[
\mathbb{D} : \mathscr{E}_{\mathrm{RSVP}} \simeq \mathscr{E}_{\mathrm{Polyx}}.
\]
\end{proposition}

This identifies geometric meaning and algebraic meaning internally.


